\section{How the Base Configuration Was Found}\label{sec:search}

This section describes the computer search that produced the base
configuration of Theorem~\ref{thm:main-seed}. The proof does not depend on
this material; we include it to document the methodology.
For the computer-assisted checks used in the proof,
see Section~\ref{sec:certified-package}.

The search requires two stages.
First, we enumerated pseudoline arrangements by generating
allowable sequences of permutations realizing the target triangle count
given by the upper bound~\eqref{eq:bbl-upper-bound}.
Second, we extracted the corresponding $\mathrm{O}$-matrices and attempted to
find straight-line realizations compatible with the constraints imposed by the
iterative constructions of~\cite{forge1998,bartholdi2007}.
The enumeration at $n=19$ was complete, but completeness is not
required: a single suitable candidate suffices for the proof of
Theorem~\ref{thm:main-seed}.

In general, stretching a pseudoline arrangement means realizing the
crossing orders encoded by its $\mathrm{O}$-matrix. If, in a fixed row, the
label $j$ must precede the label $k$, then the $y$-coordinates~\eqref{eq:y-coordinate}
of the corresponding intersection points
must satisfy an inequality of the form
\begin{equation*}
\frac{a_i-a_j}{s_i-s_j} < \frac{a_i-a_k}{s_i-s_k},
\end{equation*}
or, equivalently, after clearing denominators and taking signs into account,
in a form more suitable for computation,
\begin{equation}
(s_i-s_k)(a_i-a_j) \lessgtr (s_i-s_j)(a_i-a_k),\label{eq:general-order-ineq}
\end{equation}
where $s_i=1/m_i$. Thus the combinatorial realization problem becomes a finite
feasibility problem for row-order inequalities together with the
slope-order conditions~\eqref{eq:slope-separation}.

The present search of base configurations for the iterative construction
in Proposition~\ref{prop:iterative-step} is more restrictive.
It also fixes the $x$-intercept parameters $a_i$, while the reciprocal slope variables
$s_i=1/m_i$ remain unknown. With this choice, every adjacent row-order
comparison becomes linear in the variables $s_i$. The stretching step therefore
reduces to finding a feasible point of the system of linear inequalities
\begin{equation}
\ell_r(s_1,\dots,s_{n-1})\le -\eta,\qquad r=1,\dots,R,\label{eq:linear-order-ineq}
\end{equation}
where the linear forms $\ell_r$ are obtained from~\eqref{eq:general-order-ineq} and
$\eta>0$ is a fixed safety margin. We solve this feasibility problem with the
HiGHS linear optimization package~\cite{huangfu2018dual}, which provides
reliable infeasibility certificates.

This search generated candidate arrangements but was not part of the proof.
After obtaining a floating-point realization, the
normalization script described in Section~\ref{sec:certified-package}
determined exact rational parameters and the admissible intervals recorded in
Theorem~\ref{thm:main-seed}.
